\caption{\textbf{Statistical comparison of PhosphoFill with existing phosphorylation grafting methods.} All comparative analyses use the same 785 successfully processed single-site phosphosites. \textbf{(a,b)} Paired site-level comparison of PhosphoFill top-1 RMSD with PyTMs optimised and PTM-Psi, respectively. PhosphoFill produces the lower-RMSD placement for 619/785 sites (79\%) relative to PyTMs optimised and 670/785 sites (85\%) relative to PTM-Psi. \textbf{(c,d)} McNemar contingency analyses at a 1.0~\AA\ recovery threshold. PhosphoFill uniquely recovers 336 sites compared with 49 uniquely recovered by PyTMs optimised ($p=4.0\times10^{-48}$), and uniquely recovers 511 sites compared with 30 uniquely recovered by PTM-Psi ($p=1.3\times10^{-94}$). \textbf{(e)} Median RMSD and bootstrap 95\% confidence intervals by residue type and method, calculated using 10,000 bootstrap resamples. \textbf{(f)} Association between the PhosphoFill scoring gap and final top-1 RMSD (Spearman $\rho=0.85$, $p=1.5\times10^{-223}$). \textbf{(g)} Final top-1 RMSD stratified by the structure-specific experimental phosphoresidue environment; 773 paired sites were matched to the environment audit, with salt-bridge-like and no-contact classes shown. \textbf{(h)} Overall median RMSD across the paired benchmark cohort for all six method configurations.}
